% Generated draft: requires mathematical and source review.
\documentclass[11pt]{article}
\usepackage[margin=25mm]{geometry}
\usepackage{fontspec}
\setmainfont{Latin Modern Roman}
\usepackage{amsmath,amssymb,mathtools,unicode-math}
\setmathfont{Latin Modern Math}
\usepackage{enumitem,needspace,xurl,hyperref}
\setlength{\parindent}{0pt}
\setlength{\parskip}{5pt}
\setlength{\emergencystretch}{3em}
\newcommand{\N}{\mathbb{N}}
\newcommand{\Z}{\mathbb{Z}}
\newcommand{\Q}{\mathbb{Q}}
\newcommand{\R}{\mathbb{R}}
\newcommand{\C}{\mathbb{C}}
\newcommand{\F}{\mathbb{F}}
\newcommand{\A}{\mathbb{A}}
\newcommand{\PP}{\mathbb{P}}
\newcommand{\E}{\mathbb{E}}
\newcommand{\eps}{\varepsilon}
\newcommand{\dd}{\,\mathrm{d}}
\begin{document}
\section*{UnsolvedMath ID: 30006346}
\textbf{Source problem code:} catalogue entry
\subsection{Definitions and mathematical statement}
The supplied record concerns an array \(\Gamma\) representing a tensor and a transformed array written in the extraction as \(\Gamma X\). Its intended question is whether one can detect the existence of a prescribed null pattern and find a change of basis revealing that pattern. Schematically, the requested decision and construction task is to determine whether
\[
\exists X\;\bigl[\text{\(X\) is an admissible change of basis and the transformed array exhibits the prescribed null pattern}\bigr],
\]
and, if so, produce such an \(X\).

This is only a schematic formulation. The supplied excerpt does not define the coefficient field or ring, tensor valence, factor spaces, array dimensions, admissible changes of basis, transformation rule, or prescribed parameterized null patterns. In particular, the notation \(\Gamma X\) cannot reliably be interpreted as ordinary matrix multiplication. The quantifier above therefore has no specified domain, and the null-pattern condition has no supplied mathematical definition. These omissions prevent a precise standalone algorithmic problem statement.

There is also a substantive source-status conflict. The registry labels the item open, and its supplied literature assessment describes higher-valence tensor null-pattern detection as unresolved, adding that matrix normal forms do not extend to a general effective invariant-based solution. However, the supplied statement-verification note says that the official report already presents an algorithm which either finds a change of basis revealing the prescribed parameterized null pattern or determines that none exists. That note identifies the extraction as a reported result rather than an open problem, marks the statement unrecoverable, and says that no separate clean statement is promoted while the original record is retained publicly for provenance. The supplied material does not establish whether the apparently conflicting descriptions concern different classes of patterns or transformations. Accordingly, no precise unresolved conjecture or verified algorithmic theorem is asserted here.
\subsection{Short English statement}
The record asks about detecting tensor zero patterns and finding a basis change that reveals them, but it omits the definitions needed for a precise statement. Its open label and dated literature assessment conflict with a supplied verification note describing an existing algorithm. This statement-only draft preserves that conflict for human mathematical review.
\subsection{Sources}
\begin{itemize}
\item Supplied local registry record: permanent UnsolvedMath ID \(30006346\), source problem code \texttt{catalogue entry}, title \emph{Reported Tensor Null-Pattern Algorithm}, in the Oberwolfach Reports Open Problems collection. The supplied record has status \texttt{open}, statement status \texttt{unrecoverable}, and no clean statement.
\item Supplied report citation: \emph{Computational Group Theory} (2025), Oberwolfach Reports, pp.\ 1389--1450, DOI \texttt{10.4171/owr/2025/27}. \url{https://doi.org/10.4171/owr/2025/27}. The collection describes these reports as documenting mathematical workshops and their open-problem discussions.
\item Supplied attribution: Peter Brooksbank, \emph{Detecting structure in tensors using algebraic invariants}, joint work with Martin Kassabov and James Wilson. Supplied literature-source label: Brooksbank--Kassabov--Wilson report. \url{https://ems.press/content/serial-article-files/51860}.
\item Literature assessment checked on \texttt{2026-08-22}: the supplied curation labels the problem open and describes higher-valence tensor null-pattern detection as unresolved, with no general effective invariant-based extension of matrix normal forms. This is supplied dated curation derived from the OWR record, not an independently verified assessment; time-sensitive status and cited claims require independent verification.
\item Supplied statement-verification note: the official report is said to present an algorithm that finds a basis change revealing the prescribed parameterized null pattern or determines that none exists. The note consequently treats the extraction as a reported result rather than an open problem and retains the original record for provenance. This source-verification claim has not been independently checked here.
\end{itemize}
\end{document}
